% 第 2 章 系统总体设计 (目标 4 页)
% 对应官方要求: 技术路线(总体部分)
% 借鉴 SUFE-SAADS: 系统定位一句话 + 模块表格化 + 分层架构图 + 部署形态
% 注意: 全文自主口径, 不提任何外部项目
\section{系统总体设计}

\subsection{系统定位与设计目标}

本系统是一套部署于大模型智能体与模型服务之间的安全防护代理,对智能体的全部模型调用做透明拦截检查,覆盖攻击识别、工具调用约束、供应链风险评估、审计溯源四类能力,为政务办公场景下的智能体提供开箱可用的安全防护。

设计目标四条:

\begin{enumerate}[nosep]
  \item 高拦截。对去标记化业务话术诱导与过度代理两类高发威胁,拦截率不低于 95\%。
  \item 零误伤。正常政务办公操作的误报率不高于 1\%。
  \item 可解释。每一次拦截给出规则编号或评分理由,支持人工复核。
  \item 可复现。全部评测样本、脚本、审计日志随代码仓库提供,第三方可独立重跑。
\end{enumerate}

\subsection{总体架构}

系统采用四层架构,自上而下为接入层、策略层、判定层、审计层。

% TODO: 图 2 总体架构图(用 drawio 或 tikz 画: OpenClaw → 接入层(proxy) → 策略层(14条规则) → 判定层(LLM评审) → 审计层(audit_log), 旁路标注模型上游)
\placeholderfig{系统总体架构图(四层: 接入/策略/判定/审计)}{系统总体架构}

\begin{itemize}[nosep]
  \item \textbf{接入层}。以 HTTP 代理形态部署,实现 OpenAI Chat Completions 协议,对智能体完全透明,智能体侧仅需修改模型服务地址一项配置。
  \item \textbf{策略层}。结构化规则引擎,14 条规则覆盖已知高风险动作模式,执行开销在毫秒级。
  \item \textbf{判定层}。结构化 LLM 评审,对规则未覆盖的调用做语义级三维度评分,异常时一律阻断。
  \item \textbf{审计层}。全量调用轨迹与判定结果落盘,支持攻击复现、问题定位、效果验证。
\end{itemize}

\subsection{功能模块与赛题方向对位}

系统四大功能模块与赛题四个关键技术方向一一对应,见表~\ref{tab:module-mapping}。

\begin{table}[H]
  \centering
  \caption{功能模块与赛题关键技术方向对位}
  \label{tab:module-mapping}
  \small
  \begin{tabular}{p{2.6cm} p{1.2cm} p{5.6cm} p{4.2cm}}
    \toprule
    功能模块 & 对应方向 & 具体能力 & 实现载体 \\
    \midrule
    攻击识别与风险评估 & D1 & 对用户输入、文档附件、检索结果等多源输入做提示注入、间接指令污染、越狱诱导的自动发现与分级评估 & 请求侧扫描 + 规则层 + 评审层 \\
    工具调用安全约束 & D2 & 对文件操作、系统命令、接口调用等高风险动作做细粒度权限约束、动态策略校验与异常任务链阻断 & 代理 + 14 条结构性规则 + 评审兜底 \\
    插件与 Skill 供应链检测 & D3 & 对 Skill 包做代码行为检测、依赖关系分析与安全评级 & 50 类政务 Skill 的风险标签机制 \\
    安全评测与审计溯源 & D4 & 多维度评测体系、可量化测试方法、风险指标与审计证据输出 & 66 攻击样本 + 25 白样本 + 全量审计日志 \\
    \bottomrule
  \end{tabular}
\end{table}

四类功能以防护代理为统一入口,任一功能可独立启停,便于政企单位按实际需求裁剪部署。

\subsection{部署形态与运行环境}

系统支持两种部署形态。

\textbf{本地代理形态}。代理与智能体部署于同一主机,模型服务地址指向本地回环端口。该形态适用于单机办公智能体,部署时间不超过 10 分钟,部署步骤见附录 C。

\textbf{网关形态}。代理部署于政企内网网关,多个智能体实例共享同一代理。该形态适用于部门级部署,规则与评审配置集中管理,审计日志集中留存。

运行环境要求:Python 3.9 及以上,仅依赖标准库,无需 GPU。评审层调用远端模型服务,网络连通即可。代理自身单次调用增加的延迟实测约 1.8 秒,其中评审层占主要部分;对延迟敏感场景可仅启用策略层,增加延迟在毫秒级。
